\section{Geometry and the return problem}\label{rw:geometry}

We count cycles on one transverse section in a fixed neighborhood of
the hemicycle.  The geometric construction ensures that every counted
cycle crosses this section exactly once.  We then define the return
on the inputs whose finite return paths stay in the neighborhood;
the later changes of scale estimate this same return problem.

The base field is
\begin{equation}\label{rw:base-field}
 X_0:\qquad \dot x=-y,\qquad \dot y=x(1+y).
\end{equation}
Put $z=1+y$.  On $z>0$ it has the first integral
\begin{equation}\label{rw:geometry:energy}
 H_0(x,z)=\frac{x^2}{2}+z-1-\log z.
\end{equation}
Its Hessian is $\operatorname{diag}(1,z^{-2})$, its unique critical
point is $(0,1)$, and its sublevel sets are compact in $z>0$.
Consequently every positive level is a smooth strictly convex oval.
The field is nonzero on that oval and tangent to it, so the oval is a
periodic orbit.

We use the northern embedding
\[
 \iota(x,y)=\frac{(x,y,1)}{\sqrt{1+x^2+y^2}}.
\]
The boundary of this period annulus on the sphere is the Jordan curve
$\Gamma$ consisting of the closure of $y=-1$ and the upper equatorial
semicircle.  Indeed, at bounded $x$ the lower branch of a level tending
to infinity approaches $y=-1$.  Every unbounded sequence on these
ovals has $y>-1$ and hence can accumulate only on the upper semicircle;
intersections with each fixed ray $y=k|x|$, $k\ge0$, tend to infinity,
as does the upper turning point.  Write $p_-,p_+$ for the negative and
positive horizontal points at infinity, and $q_+$ for the positive
vertical point.  The direction on $y=-1$ is from $p_-$ to $p_+$.
In the chart $y=1/r$, $x=w/r$, multiplication by $r>0$ gives
\[
 r'=-rw(1+r),\qquad w'=-w^2-r(1+w^2).
\]
Thus the other half of $\Gamma$ runs from $p_+$ through $q_+$ to $p_-$.

\subsection{A coefficient family for all quadratic perturbations}

Consider the five-parameter family
\begin{equation}\label{rw:family}
 \begin{aligned}
 X_\alpha=(P_\alpha,Q_\alpha),\qquad
 P_\alpha&=-y+Bx^2+my^2+ax,\\
 Q_\alpha&=x(1+y)+cy+dy^2,
 \qquad \alpha=(B,m,a,c,d).
 \end{aligned}
\end{equation}
The coefficient which breaks the invariant line is
\begin{equation}\label{rw:geometry:line-breaking}
 \ell=d-c,\qquad \dot z\big|_{z=0}=\ell.
\end{equation}
The origin remains an equilibrium; its linearization has determinant
$1+ac>0$ near $\alpha=0$.

\begin{proposition}\label{rw:slice}
The family \eqref{rw:family}, near-identity affine coordinate changes,
and positive constant time changes form analytic local coordinates on
the twelve-dimensional space of real quadratic planar vector fields
near $X_0$.  A uniform bound for isolated cycles of \eqref{rw:family}
in a fixed spherical collar of $\Gamma$, counted with return-germ
multiplicity, therefore gives the same bound on a smaller fixed collar
for every quadratic field in an open coefficient neighborhood of $X_0$.
\end{proposition}

\begin{proof}
For $A$ near the identity matrix, $b\in\R^2$ near zero, and $r\in\R$
near zero, define
\[
 \mathcal S(A,b,r,\alpha)(v)
       =e^r A^{-1}X_\alpha(Av+b).
\]
There are four matrix entries, two translation parameters, one time
parameter, and five parameters in $\alpha$.
For an affine infinitesimal change $h$, its column in the differential
at the base point is $DX_0h-DhX_0$.  The changes
$(1,0),(0,1),(x,0),(y,0),(0,x),(0,y)$ give, respectively,
\[
 \begin{array}{lll}
 v_1=(0,1+y),&v_2=(-1,x),&v_3=(y,x+xy),\\
 v_4=(-x-xy,y+y^2),&v_5=(-x,x^2+y),&v_6=(-y,-x).
 \end{array}
\]
The time column is $v_7=(-y,x+xy)$.  The remaining columns, in the
parameter order of \eqref{rw:family}, are
\[
 (x^2,0),\quad(y^2,0),\quad(x,0),\quad(0,y),\quad(0,y^2).
\]
In a linear relation, the constant, first-component $xy$ and
second-component $x^2$ coefficients eliminate $v_1,v_2,v_4,v_5$.
The remaining coefficients $u,v,w$ of $v_3,v_6,v_7$ satisfy
\[
 \begin{pmatrix}1&-1&-1\\1&-1&1\\1&0&1\end{pmatrix}
 \begin{pmatrix}u\\v\\w\end{pmatrix}=0.
\]
Its determinant is $-2$, and the five family columns occupy distinct
monomials. Thus $D\mathcal S$ is invertible; the analytic inverse
function theorem proves the local-coordinate assertion.

If $Y=\mathcal S(A,b,r,\alpha)$ and $\Phi(v)=Av+b$, then
$D\Phi\,Y=e^r X_\alpha\circ\Phi$.  Thus $\Phi$ sends trajectories to
trajectories preserving time orientation.  On transverse sections it
conjugates the corresponding return germs.  If $\psi$ is the section
coordinate change and $P$ a return, then
\begin{equation}\label{rw:geometry:section-unit}
 \psi(P(u))-\psi(u)
 =(P(u)-u)\int_0^1\psi'\bigl(u+t(P(u)-u)\bigr)\,\dd t.
\end{equation}
At a fixed point the last factor is nonzero.  Isolation and the order
of the displacement zero are preserved.

In homogeneous coordinates the spherical extension of $\Phi$ is
\[
 \widehat\Phi(\xi_1,\xi_2,\xi_3)
 =\frac{(A(\xi_1,\xi_2)^{\mathsf T}+b\xi_3,\xi_3)}
 {\left\|(A(\xi_1,\xi_2)^{\mathsf T}+b\xi_3,\xi_3)\right\|}.
\]
Its homogeneous matrix is invertible, so it is a sphere
diffeomorphism preserving the closed northern hemisphere.  It converges
uniformly to the identity as $(A,b)$ tends to $(I,0)$.  Given a collar
$U_*$ for the normalized family, choose $U$ with
$\overline U\subset U_*$.  One sufficiently small affine neighborhood
then satisfies $\widehat\Phi(\overline U)\subset U_*$ for every such
$\Phi$.  Intersecting the corresponding product neighborhood with the
local-coordinate neighborhood of $\mathcal S$ gives an open
twelve-dimensional coefficient neighborhood.  Cycles in $U$ map
injectively to cycles in $U_*$ with their multiplicities unchanged.
\end{proof}

\subsection{A fixed collar and its transverse section}

The initial energy count uses the base-field Chebyshev estimate of
Proposition~\ref{rw:base-ect}.  Fix its threshold
$H_{\mathrm{ECT}}$ and choose the collar so that every point on its
positive lower section has energy above that threshold.

\begin{proposition}\label{rw:collar}
There are a closed spherical annulus $\mathcal A$ containing $\Gamma$
in its interior, a number $0<z_L<1/8$, and a coefficient neighborhood
$\mathcal V$ of zero with the following properties.  The annulus
excludes the origin and the negative vertical point at infinity, and
\begin{equation}\label{rw:geometry:section}
 h(z_L):=z_L-1-\log z_L>H_{\mathrm{ECT}}+2,
 \qquad \Sigma=\{x=0,\ |z|<z_L\}.
\end{equation}
Its lower part over $-1\le x\le1$ is the rectangle
$[-1,1]\times[-z_L,z_L]$ in $(x,z)$ coordinates.
Every affine periodic orbit contained in $U=\operatorname{int}\mathcal A$
is essential in $\mathcal A$ and crosses $\Sigma$ exactly once, with
$\dot x>0$.  Both $\mathcal A$ and $\mathcal V$ can be chosen invariant
under the spatial and coefficient reflections defined below.
\end{proposition}

\begin{proof}
We first choose three vertex boxes in which no periodic orbit can be
contained.  Throughout this construction let all five coefficient
moduli be at most $p_*=1/64$.
At a horizontal vertex put $x=\sigma/r$, $y=v/r$, where
$\sigma\in\{-1,1\}$, and multiply the field by $r>0$.  Its equations
are
\[
 \begin{aligned}
 r'&=-\sigma Br-\sigma mrv^2+\sigma r^2v-ar^2,\\
 v'&=\sigma r+\sigma(1-B)v+dv^2-\sigma mv^3
          +\sigma rv^2+(c-a)rv.
 \end{aligned}
\]
The divergence of this polynomial field is
\[
 \sigma(1-2B)+2dv-4\sigma mv^2+4\sigma rv+(c-3a)r.
\]
On the square $|r|,|v|\le1/8$ its difference from $\sigma$ is at
most $109/1024<1/8$.  If a periodic orbit were contained in its
northern part $r>0$, its Jordan disk would lie in the same square
and half-plane.  Green's theorem would give zero for the integral
of the divergence over that disk, contrary to its strict sign.

At the upper vertex use $y=1/r$, $x=w/r$.  The positively rescaled
field is $r'=-rF$, $w'=G$, where
\[
 \begin{aligned}
 F&=(1+r)w+d+rc,\\
 G&=m-dw+(B-1)w^2+r[-1+(a-c)w-w^2].
 \end{aligned}
\]
In $0<r<1/4$, $|w|<1/4$, an equilibrium lies on
\[
 w=w_s(r)=-\frac{d+rc}{1+r},\qquad
 |w_s|\le p_*,\quad |w_s'|\le2p_*.
\]
On this graph,
\[
 G_r\le-1+2p_*^2,\qquad |G_w|\le4p_*+2p_*^2,
\]
so
\[
 \frac{\dd}{\dd r}G(r,w_s(r))
 \le-1+10p_*^2+4p_*^3<-\frac12.
\]
There is at most one equilibrium in the box, and its determinant is
\[
 \det D(-rF,G)
 =r(1+r)\frac{\dd}{\dd r}G(r,w_s(r))<0.
\]
It is a saddle.  A periodic orbit contained in this rectangle would
enclose either no equilibrium or one saddle, whereas its index is one.
This excludes such an orbit.

Cut the regular parts of $\Gamma$ at finitely many points, with two
cuts inside each vertex box.  The remaining compact regular arcs have
flow-box neighborhoods.  A sufficiently thin bicollar can be divided
by radial crosscuts into strips, each contained in one of these flow
boxes or vertex boxes.  Make the crosscuts transverse at the chosen
regular points, with their positive normal pointing in the direction
of the base flow.  Include the lines $x=-1,0,1$ on the finite arc
and choose the stated rectangular part there.  Its width can be made
arbitrarily small; choose $z_L$ to satisfy
\eqref{rw:geometry:section}.  The other widths can then be decreased
so that the annulus excludes the two points specified in the statement.
The construction can be made symmetric under $x\mapsto-x$ by choosing
the two horizontal boxes and the regular crosscuts in reflected pairs.

On the central rectangle, $P_\alpha>3/4$ after shrinking the
coefficient neighborhood.  On every other crosscut the normal
component also has a positive minimum on its compact length.  The
finitely many strict inequalities persist on one common coefficient
neighborhood, which we intersect with its reflected image.

Any periodic orbit in $U$ must meet a crosscut.  Otherwise it is
contained in a single strip, where a flow-box coordinate or one of the
vertex arguments excludes it.  All crosscut intersections are transverse
and positive.  There are finitely many, since the orbit is compact and
a transverse intersection is isolated.  Each crosscut represents the
same relative homology class in the annulus, so its signed intersection
number with the orbit is the orbit's winding number.  A simple closed
curve in an annulus has winding number $0$ or $\pm1$; positivity
forces the value one.  Hence the orbit is essential and meets every
crosscut, in particular $\Sigma$, exactly once.
\end{proof}

\subsection{Actual returns and periodic points}

Only isolated affine periodic orbits whose northern images lie in $U$
are counted.  We first locate these orbits relative to the line $z=0$
and bound the line-breaking coefficient by their section height.  The
return domain will then be defined by the complete trajectory.

\begin{lemma}\label{rw:geometry:cycle-height}
Every periodic orbit contained in $U$ encloses the origin and lies in
$z>0$.  If $s$ is its intersection with $\Sigma$, then
\begin{equation}\label{rw:geometry:line-bound}
 0<s<z_L,\qquad |\ell|<3s.
\end{equation}
\end{lemma}

\begin{proof}
An essential curve in $\mathcal A$ separates its two complementary
components.  The origin is on the bounded side of $\Gamma$, while
the southern hemisphere is on the other side.  Thus the affine Jordan
disk of each such periodic orbit contains the origin.

For $\ell\ne0$, every intersection with $z=0$ has the same nonzero
normal velocity $\ell$.  A periodic scalar function cannot have only
crossings of one sign, so a periodic orbit cannot meet this line.
For $\ell=0$, invariance and uniqueness give the same exclusion.
Choose one of the regular equatorial crosscuts of the preceding
construction with its northern part in $y>0$.  Every counted orbit
meets it, and hence is on the side $z>0$ throughout.

Its lower segment is a graph over $[-1,1]$, since $P_\alpha>3/4$
on the central rectangle.  On that graph the exact equation is
\[
 z_x=\ell A(x,z)+b(x,z)z,
 \quad A=\frac{(1-z)^2}{D},\quad
 b=\frac{x+c(z-1)}{D},
\]
where
\[
 D=1-z+Bx^2+m(1-z)^2+ax.
\]
For $0\le z\le z_L<1/8$ and the coefficient bound used above,
$3/4<D<9/8$, $A>2/3$, and $|b|<3/2$.
If $\ell<0$, integrate from $0$ to $1$ and use $z(1)>0$ to obtain
\[
 s>|\ell|\int_0^1 A(t,z(t))
             e^{-\int_0^t b(u,z(u))\,\dd u}\,\dd t
 \ge\frac49(1-e^{-3/2})|\ell|>\frac{|\ell|}{3}.
\]
If $\ell>0$, variation of constants from $-1$ to $0$ gives the same
lower bound for $s$ after discarding its positive homogeneous term.
The case $\ell=0$ is immediate.
\end{proof}

Let $\Sigma_+=\{x=0,\ 0<z<z_L\}$, identified with $(0,z_L)$.
For a fixed $\alpha$, define $\mathcal D_\alpha$ to be the set of
$s\in\Sigma_+$ whose trajectory has a finite first positive-time
return to $\Sigma_+$, with the entire intervening path affine,
contained in $U$, and contained in $z>0$.  Write this return as
$\mathcal P_\alpha(s)$.

\begin{lemma}\label{rw:geometry:return-domains}
The domain $\mathcal D_\alpha$ is open.  The return is locally jointly
analytic in its input and coefficients and is a local increasing
diffeomorphism.  Its isolated fixed points correspond bijectively to
the cycles counted in $U$.
\end{lemma}

\begin{proof}
A finite return path is compact in the prescribed open region and
ends transversely. Analytic flow dependence and the implicit function
theorem preserve its return under small input and coefficient changes.
Compactness away from the endpoint flow boxes, and transversality
inside them, exclude an earlier intersection. This proves openness
and local joint analyticity.
The variational determinant formula for a transverse return gives
\[
 \mathcal P_\alpha'(s)=
 \frac{P_\alpha(0,s-1)}
      {P_\alpha(0,\mathcal P_\alpha(s)-1)}
 \exp\!\left(\int_0^{T(s)}\operatorname{div}X_\alpha
                          (\Phi_t^\alpha(0,s-1))\,\dd t\right)>0.
\]
Here $P_\alpha$ is the first component of the vector field in
\eqref{rw:family}.
Both normal velocities are positive.  Finally, a fixed point closes
the prescribed trajectory, and a counted cycle has exactly one section
intersection by Proposition~\ref{rw:collar} and belongs to this domain
by Lemma~\ref{rw:geometry:cycle-height}.
\end{proof}

\begin{definition}\label{rw:geometry:counting}
The Poincar\'e multiplicity of an isolated cycle is the order of the
zero of $\mathcal P_\alpha(s)-s$ at its section intersection.  A return germ
identically equal to the identity contributes no isolated cycle.
\end{definition}

This definition is independent of the analytic section coordinate by
\eqref{rw:geometry:section-unit}.  Multiplication of a displacement by
an analytic nonvanishing factor also preserves its zero orders.
Derivative zeros require a separate comparison, since differentiating
a unit multiple also differentiates the unit.  The domain
$\mathcal D_\alpha$ may be disconnected.  Thus a periodic orbit
supplies a local return, while a comparison with more distant section
points requires the intervening paths to be constructed.

\subsection{One direction for each fixed field}

The reflection used throughout the paper reverses time as well as space.
Let
\begin{equation}\label{rw:geometry:reflection}
 R(x,y)=(-x,y),\qquad
 \alpha^R=(B,m,-a,-c,-d).
\end{equation}
Direct substitution in \eqref{rw:family} gives
\begin{equation}\label{rw:geometry:flow-reflection}
 X_{\alpha^R}(R v)=-DR\,X_\alpha(v),\qquad
 \Phi_t^{\alpha^R}(Rv)=R\Phi_{-t}^{\alpha}(v).
\end{equation}
The collar and coefficient neighborhood are reflection invariant, and
$R$ fixes every point of $\Sigma_+$.  Thus it preserves cycles, their
central height, and isolation.

Let $s_*$ be a periodic point and restrict its original
forward return to an analytic increasing diffeomorphism
$P:I\to P(I)$ on a small positive section interval.  The reflected
forward return near that point is $Q=P^{-1}:P(I)\to I$, and
\begin{equation}\label{rw:geometry:inverse-germ}
 Q(s)-s=-\bigl(P(Q(s))-Q(s)\bigr).
\end{equation}
Since $Q'(s_*)>0$, composition with $Q$ preserves the order of the
displacement zero.  The multiplicity is therefore unchanged by
\eqref{rw:geometry:reflection}.  This is an identity of germs at a
periodic point.  The reflected positive return follows the original
trajectory in reverse time, so its domain is determined by that reverse
passage.

We may consequently reflect a fixed field with $\ell<0$ and carry out
its count with $\ell\ge0$.  All section labels and comparison paths
for that field are thereafter taken in this selected direction.
Numerical cutoffs used uniformly in the family are chosen once on a
reflection-invariant coefficient neighborhood, not chosen anew for
each field.  This sign-dependent choice is made only after the field
is fixed.  Coefficient differentiation and division across $\ell=0$
use the original analytic auxiliary maps on their stated common domains.

\subsection{Energy ranges and the primary normalization}

The lower section has the analytic coordinate
\[
 h=h(s)=s-1-\log s>H_{\mathrm{ECT}}+2,
 \qquad h'(s)=1-s^{-1}<0.
\]
For $0<\eta<1$ and a fixed $H_*>h(z_L)+1$, set
\[
 \mathscr M(h)=\max\{|B|h,|m|h,|a|\sqrt h,|c|\sqrt h\}.
\]
The three energy classes are
\begin{equation}\label{rw:geometry:energy-classes}
 \begin{array}{ll}
 h\le H_*,&\text{bounded energy},\\
 h>H_*,\quad\mathscr M(h)\le\eta,
       &\text{small coefficients at the energy scale},\\
 h>H_*,\quad\mathscr M(h)>\eta,
       &\text{parameter-dominant energy}.
 \end{array}
\end{equation}
The first two classes will be counted together.  Their thresholds are
fixed before the final coefficient shrink, and the bounded-energy
interval lies entirely above $H_{\mathrm{ECT}}$.
This classification applies to actual periodic
points and extends no return domain.

\begin{definition}\label{rw:normalization}
The primary scale and normalized coefficients, for $\rho>0$, are
\begin{equation}\label{rw:geometry:primary-scale}
 \begin{gathered}
 \rho=\max\{\sqrt{|B|},\sqrt{|m|},|a|,|c|\},
 \qquad\epsilon=\rho^2,\\
 (B,m,a,c)=(\rho^2\bar B,\rho^2\bar m,
                       \rho\bar a,\rho\bar c),\qquad
 \nu=\ell/\rho,\quad\bar d=d/\rho=\bar c+\nu.
 \end{gathered}
\end{equation}
Thus $\max(|\bar B|,|\bar m|,|\bar a|,|\bar c|)=1$.
The generators in original and primary coefficients are
\begin{equation}\label{rw:geometry:generators}
 \begin{gathered}
 \tau_0=a+c,\qquad\sigma_0=c(B+m),\\
 \tau_1=\bar a+\bar c=\tau_0/\rho,\qquad
 \sigma_1=\bar c(\bar B+\bar m)=\sigma_0/\rho^3.
 \end{gathered}
\end{equation}
\end{definition}

The normalized vector lies on the union of the signed faces of
$[-1,1]^4$.  Smooth parameter arguments are made in the indicated
local coefficient charts, not by differentiating the maximum across
its corners.  For $\rho>0$, the two products $\tau_0\sigma_0$ and
$\tau_1\sigma_1$ have the same sign, but
\begin{equation}\label{rw:geometry:generator-ratio}
 \frac{\tau_1}{\sigma_1}
       =\epsilon\frac{\tau_0}{\sigma_0}\qquad(\sigma_0\ne0).
\end{equation}
If $\rho=0$, then $B=m=a=c=0$, while $d=\ell$ may be nonzero.
All energies belong to the first two classes because $\mathscr M=0$.
No barred coefficients, $\nu$, or quotient by $\epsilon$ is used in
that case.

For later calculations we record the exact scaled equations.  With
$\xi=t/\rho$, $X=\rho x$, and $z=1+y$, the lower equations are
\begin{equation}\label{rw:geometry:lower-field}
 \begin{aligned}
 X_\xi&=\epsilon\{1-z+\bar B X^2+\bar aX
                              +\epsilon\bar m(1-z)^2\},\\
 z_\xi&=Xz+\epsilon\bar c\,z(z-1)+\rho\ell(1-z)^2.
 \end{aligned}
\end{equation}
With $X=\rho x$ and $Y=\rho^2y$, the same field and time give
\begin{equation}\label{rw:geometry:upper-field}
 \begin{aligned}
 X_\xi&=Y(\bar mY-1)+\epsilon(\bar B X^2+\bar aX),\\
 Y_\xi&=(Y+\epsilon)(X+\bar cY)+\nu Y^2.
 \end{aligned}
\end{equation}
These identities do not assume that $\nu$ is bounded.

\begin{lemma}\label{rw:geometry:flatness}
A parameter-dominant cycle has $\rho>0$ and satisfies
\[
 \rho\sqrt h>\eta,\qquad
 |\ell|<3e^{-\eta^2/\rho^2},\qquad
 \frac{|\ell|}{\rho^k}\le C_{k,\eta}
                  e^{-\eta^2/(2\rho^2)}\quad(k\ge0).
\]
\end{lemma}
\begin{proof}
If $\rho\sqrt h\le\eta$, the four terms in $\mathscr M(h)$ are
bounded by $\eta^2,\eta^2,\eta,\eta$, respectively, contradicting
the third class.  Since $0<s<1$, the exact identity
$s=e^{-h-1+s}$ gives $s<e^{-h}$.  Apply
\eqref{rw:geometry:line-bound} and retain half of the exponential to
absorb any fixed power of $\rho^{-1}$.
\end{proof}
The estimate uses the height of an actual cycle.  The ratio $\ell/s$
still has only the bound in \eqref{rw:geometry:line-bound}; derivatives
along parameter curves require the common-domain estimates below.

\begin{lemma}\label{rw:geometry:endpoints}
Every periodic orbit counted in $U$ meets $y=0$ exactly once on each
half-axis.  For $\rho>0$, write these intersections, in the selected
direction, as $(p/\rho,0)$ and $(-q/\rho,0)$ with $p,q>0$.
The upper connecting arc lies in $y>0$.  Nested cycles have increasing
$p$ and $q$ towards the outer cycle.  Reflection exchanges the labels:
$p_{\rm ref}=q_{\rm orig}$ and $q_{\rm ref}=p_{\rm orig}$.
\end{lemma}
\begin{proof}
By Lemma~\ref{rw:geometry:cycle-height}, the Jordan disk contains
the origin.  Each horizontal half-axis therefore meets its boundary.
On $y=0$, the normal velocity is $\dot y=x$, so all crossings of
the positive ray are upward and all crossings of the negative ray
are downward.  Their signed intersection number is the winding number
one; consequently each ray is crossed exactly once.  The arc between
the upward and downward crossings is in $y>0$.  Inclusion of the
Jordan disks orders their two ray intersections and proves the nesting
assertion.  Equation~\eqref{rw:geometry:flow-reflection} exchanges the
two crossings while reversing time, giving the last assertion.
\end{proof}

In particular a cutoff imposed on the reflected outgoing label is a
cutoff on the original incoming label, not on the original outgoing
label.  All later endpoint cutoffs refer to the direction already
selected for the fixed field.  Varying a section input always keeps
the original coefficient vector fixed; any induced changes of scaled
coordinates must be included in the chain rule.
